\Needspace{9\baselineskip}\begin{center}\itshape \textit{Reducing the equations used to find the roots \(\Omega\) to pure power equations.}\end{center}


\Needspace{6\baselineskip}\Article{359}
The preceding investigations concerned the \textit{finding} of auxiliary equations; we will now explain a particularly remarkable property concerning their \textit{solution}. As is well known, all the efforts of the greatest geometers to find the general solution of equations beyond the fourth degree, or (to define more precisely what I am seeking) \textsc{the reduction of mixed equations to pure power equations}, have so far always been fruitless, and scarcely any doubt remains that this problem proposes something impossible rather than something beyond the powers of present-day analysis. (Compare our remarks on this subject in \textit{Demonstr. nova etc. art. 9}.) Nevertheless, it is certain that there are countless mixed equations of every degree that admit such a reduction to pure power equations, and we hope that geometers will welcome our showing that our auxiliary equations always belong to this class. But because of the great extent of this investigation,
\SourceLeaf{455}
we give here only the principal points needed to show the possibility, and defer to another time the fuller treatment that this subject so richly deserves. We must first make some general observations about the roots of the equation \(x^e-1=0\), including the case where \(e\) is composite.

I. These roots are represented (as is known from elementary books) by \(\cos\frac{kP}e+i\sin\frac{kP}e\), where \(k\) is to take the \(e\) numbers \(0,\allowbreak 1,\allowbreak 2,\allowbreak 3\ldots e-1\), or any others congruent to these modulo \(e\). One root, for \(k=0\), or generally for \(k\) divisible by \(e\), becomes \(=1\); every other value of \(k\) corresponds to a root different from \(1\).

II. Since \((\cos\frac{kP}e+i\sin\frac{kP}e)^\lambda=\cos\frac{\lambda kP}e+i\sin\frac{\lambda kP}e\), it is clear that if \(R\) is a root corresponding to a value of \(k\) coprime to \(e\), in the progression \(R,\allowbreak RR,\allowbreak R^3\), etc., the term in position \(e\) is indeed \(=1\), but all the preceding terms differ from \(1\). It follows at once that all \(e\) quantities \(1,\allowbreak R,\allowbreak RR,\allowbreak R^3\ldots R^{e-1}\) are unequal, and since they evidently all satisfy the equation \(x^e-1=0\), they represent all its roots.

III. Finally, under the same assumption, we have \[1+R^\lambda+R^{2\lambda}\ldots+R^{\lambda(e-1)}\ =0\] for every integer value of \(\lambda\) not divisible by \(e\); for the sum is \(=\frac{1-R^{\lambda e}}{1-R^\lambda}\), and the numerator of this fraction becomes \(=0\), whereas its denominator does not become \(=0\). When, however, \(\lambda\) is divisible by \(e\), that sum evidently becomes \(=e\).

\Needspace{6\baselineskip}\Article{360}
Let, as throughout the preceding discussion, \(n\) be a prime, \(g\) a primitive root modulo \(n\), and \(n-1\) the product of three positive integers \(\alpha,\allowbreak \beta,\allowbreak \gamma\). For brevity we will immediately set up the investigation so that it also includes the cases where \(\alpha\) is one or \(\gamma=1\); when \(\gamma=1\), the sums \((\gamma,1),\allowbreak (\gamma,g)\), etc., must be replaced by the roots \([1],\allowbreak [g]\), etc. We therefore suppose that from all the \(\alpha\) known sums of \(\beta\gamma\) terms \((\beta\gamma,1),\allowbreak (\beta\gamma,g),\allowbreak (\beta\gamma,gg)\ldots(\beta\gamma,g^{\alpha-1})\) we are to derive the sums of \(\gamma\) terms. We previously reduced this task to a mixed equation of degree \(\beta\); we will now show how to accomplish it by a pure power equation of the same degree. For brevity, we will write, for the sums \[(\gamma,1),(\gamma,g^\alpha),(\gamma,g^{2\alpha})\ldots(\gamma,g^{\alpha\beta-\alpha})\]
\SourceLeaf{456}
contained in \((\beta\gamma,1)\), \(a,\allowbreak b,\allowbreak c\ldots m\) respectively; for the sums \[(\gamma,g),(\gamma,g^{\alpha+1})\ldots(\gamma,g^{\alpha\beta-\alpha+1})\] contained in \((\beta\gamma,g)\), \(a',\allowbreak b'\ldots m'\) respectively; for the sums \[(\gamma,gg),(\gamma,g^{\alpha+2})\ldots(\gamma,g^{\alpha\beta-\alpha+2})\] we write \(a'',\allowbreak b''\ldots m''\) respectively; and so on as far as those contained in \((\beta\gamma,g^{\alpha-1})\).

I. Now let \(R\) denote an unspecified root of the equation \(x^\beta-1=0\), and suppose that expanding the power with exponent \(\beta\) of the function \[t=a+Rb+RRc\ldots+R^{\beta-1}m\] according to the rules of Article 345 gives \[\begin{aligned}&N+Aa+Bb+Cc\ldots+Mm\\&{}+A'a'+B'b'+C'c'\ldots+M'm'\\&{}+A''a''+B''b''+C''c''\ldots+M''m''\\&{}+\text{etc.}=T\end{aligned}\] where all the coefficients \(N,\allowbreak A,\allowbreak B,\allowbreak A'\), etc., will be polynomials in \(R\). Suppose also that the powers with exponent \(\beta\) of two other functions \[\begin{aligned}u&=R^\beta a+Rb+RRc\ldots+R^{\beta-1}m,\\u'&=b+Rc+RRd\ldots+R^{\beta-2}m+R^{\beta-1}a\end{aligned}\] expand into \(U\) and \(U'\) respectively. It is easy to see from Article 350 that, since \(u'\) arises from \(t\) by interchanging the sums \(a,\allowbreak b,\allowbreak c\ldots m\) with \(b,\allowbreak c,\allowbreak d\ldots a\) respectively, we will have \[\begin{aligned}&U'=N+Ab+Bc+Cd\ldots+Ma\\&{}+A'b'+B'c'+C'd'\ldots+M'a'\\&{}+A''b''+B''c''+C''d''\ldots+M''a''\\&{}+\text{etc.}\end{aligned}\]

Further, since \(u=Ru'\), we clearly have \(U=R^\beta U'\), so that, because \(R^\beta=1\), the corresponding coefficients in \(U\) and \(U'\) are equal. Finally, since \(t\) and \(u\) differ only in that \(a\) is multiplied by unity in \(t\) and in \(u\) by \(R^\beta\), it is easy to see that all corresponding coefficients (that is, those multiplying the same sums) in \(T\) and \(U\) are equal, and therefore so are all corresponding coefficients in \(T\) and \(U'\). Hence we finally conclude that \(A=B=C\), etc., \(=M\);
\SourceLeaf{457}
\(A'=B'=C'\), etc., \(A''=B''=C''\), etc., etc. Thus \(T\) reduces to a form such as \[N+A(\beta\gamma,1)+A'(\beta\gamma,g)+A''(\beta\gamma,gg)\ \text{etc.}\] where each coefficient \(N,\allowbreak A,\allowbreak A'\), etc., may be reduced to a form such as {\postdisplaypenalty=10000 \[pR^{\beta-1}+p'R^{\beta-2}+p''R^{\beta-3}+\text{etc.}\]} in which \(p,\allowbreak p',\allowbreak p''\), etc., are given integers.

II. If for \(R\) we take a definite root of the equation \(x^\beta-1=0\) (whose solution we suppose already known), and specifically a root of which no power lower than the one with exponent \(\beta\) equals unity, then \(T\) too will be a definite quantity, from which \(t\) may be derived by the pure power equation \(t^\beta-T=0\). But since this equation has \(\beta\) roots, namely \(t,\allowbreak Rt,\allowbreak RRt\ldots R^{\beta-1}t\), it may seem doubtful which root we ought to adopt. That the choice is entirely arbitrary will, however, become clear as follows. We must remember that, after all the sums of \(\beta\gamma\) terms have been determined, the root \([1]\) is determined only to the extent that some one of the \(\beta\gamma\) roots contained in \((\beta\gamma,1)\) must be denoted by this symbol; and therefore it is entirely arbitrary which of the \(\beta\) sums making up \((\beta\gamma,1)\) we choose to denote by \(a\). Now suppose that, with some definite sum denoted by \(a\), we obtain \(t=\mathfrak{T}\). It is easy to see that if we subsequently wish the sum just denoted by \(b\) to be denoted by \(a\), those previously denoted by \(c,\allowbreak d\ldots a,\allowbreak b\) now become \(b,\allowbreak c\ldots m,\allowbreak a\), and the value of \(t\) now becomes \(=\frac{\mathfrak{T}}R=\mathfrak{T}R^{\beta-1}\). Similarly, if we choose to denote by \(a\) the sum originally denoted by \(c\), the value of \(t\) will become \(\mathfrak{T}R^{\beta-2}\), and so on: \(t\) can be regarded as equal to any of the quantities \(\mathfrak{T},\allowbreak \mathfrak{T}R^{\beta-1},\allowbreak \mathfrak{T}R^{\beta-2}\), etc., \(\text{that is}\) to any root of the equation \(x^\beta-T=0\), according as we suppose one or another of the sums contained in \((\beta\gamma,1)\) to be denoted by \((\gamma,1)\). \(Q.\ E.\ D.\)

III. After the quantity \(t\) has been determined in this way, we must find \(\beta-1\) others that arise from \(t\) when, in its expression, \(R\) is successively replaced by \(RR,\allowbreak R^3,\allowbreak R^4\ldots R^\beta\), namely \[\begin{aligned}t'&=a+RRb+R^4c\ldots+R^{2\beta-2}m,\\t''&=a+R^3b+R^6c\ldots+R^{3\beta-3}m\ \text{etc.}\end{aligned}\] The last is already known, since it evidently becomes \(=a+b+c\ldots+m=(\beta\gamma,1)\); the others may be obtained as follows. If, according to the rules of Article 345 and in the same way as \(t^\beta\) earlier in I, we expand the product \(t^{\beta-2}t'\), it will be proved by a method entirely analogous to the preceding one that the result can be reduced to a form such as
\SourceLeaf{458}
\[\mathfrak{N}+\mathfrak{A}(\beta\gamma,1)+\mathfrak{A}'(\beta\gamma,g)+\mathfrak{A}''(\beta\gamma,gg)\ \text{etc.}=T'\] in which \(\mathfrak{N},\allowbreak \mathfrak{A},\allowbreak \mathfrak{A}'\), etc., are polynomials in \(R\), so that \(T'\) is a known quantity, from which we will obtain \(t'=\frac{T'tt}{T}\). In exactly the same way, if we suppose that expanding the product \(t^{\beta-3}t''\) gives \(T''\), this expression will have a similar form, and from its known value we will therefore derive \(t''\) by the equation \(t''=\frac{T''t^3}{T}\); similarly, \(t'''\) will be found by an equation such as \(t'''=\frac{T''''t^4}{T}\), in which \(T'''\) is a known quantity, and so on.

This method would not be applicable if \(t=0\) could occur, whence we would also have to have \(T=T'=T''\), etc., \(=0\); but this can be proved impossible, although we must omit the proof here because of its length. — There are also special devices for converting the fractions \(\frac{T'}T,\allowbreak \frac{T''}T\), etc., into \textit{polynomials} in \(R\), as well as shorter methods, in the case \(\alpha=1\), of obtaining the values of \(t',\allowbreak t''\), etc. We must pass over all these here without discussion.

IV. Finally, as soon as \(t,\allowbreak t',\allowbreak t''\), etc., have been found, observation III of the preceding article immediately gives \(t+t'+t''+\text{etc.}=\beta a\), so that the value of \(a\) will be known, from which the values of all the remaining sums of \(\gamma\) terms may be derived by Article 346. — The values of \(b,\allowbreak c,\allowbreak d\), etc., can also be obtained from the following equations, whose reason will be readily apparent to any attentive reader: \[\begin{aligned}\beta b&=R^{\beta-1}t+R^{\beta-2}t'+R^{\beta-3}t''+\text{etc.}\\\beta c&=R^{2\beta-2}t+R^{2\beta-4}t'+R^{2\beta-6}t''+\text{etc.}\\\beta d&=R^{3\beta-3}t+R^{3\beta-6}t'+R^{3\beta-9}t''+\text{etc. etc.}\end{aligned}\]

Of the many observations relating to the preceding investigation, we touch on only one here. As regards the solution of the pure power equation \(x^\beta-T=0\), it is clear that \(T\) in most cases has an imaginary value \(P+iQ\), so that this solution will, as is known, depend partly on dividing an angle (whose tangent \(=\frac QP\)) and partly on dividing a ratio (unity to \(\surd(PP+QQ)\)) into \(\beta\) parts. It is very remarkable here (though we do not pursue this in greater detail) that the value of \({}^{\beta}\!\surd(PP+QQ)\) can always be expressed \textit{rationally} in terms of quantities already known, so that, apart from extracting a square root, \textit{only} the division of an angle is required for the solution, \(\text{e.g.}\) for \(\beta=3\) only the trisection of an angle.

Finally, since nothing prevents our setting \(\alpha=1,\allowbreak \gamma=1\) and therefore
\SourceLeaf{459}
\(\beta=n-1\), it is evident that solving the equation \(x^n-1=0\) can at once be reduced to solving the pure power equation of degree \(n-1\), namely \(x^{n-1}-T=0\), where \(T\) will be determined by the roots of the equation \(x^{n-1}-1=0\). From this, with the help of the observation just made, it follows that dividing the whole circle into \(n\) parts requires \(1^\circ\) dividing the whole circle into \(n-1\) parts, \(2^\circ\) dividing another arc, which can be constructed after that division has been made, into \(n-1\) parts, \(3^\circ\) extracting a single square root; and it can indeed be shown that this is always \(\surd n\).
